\chapter{External Materials}
\label{app:external}

This appendix records what external material accompanies the dissertation, where
each item sits, how the study is regenerated from scratch, and how any number
printed in the preceding chapters can be traced back to the artefact that
produced it. This report is submitted on its own, and everything behind it is shared through a
public repository at \url{https://github.com/Beiciccc/sp500vol-bench}, open to the
supervisor, the assessors and any other reader without an account or a request.
Every artefact described in Sections~\ref{app:external:release}
to~\ref{app:external:process} is in it, and none of those sections depends on a
resource a marker cannot open. What the repository does not carry is what the
licence bars from redistribution --- the returns, the labels, the price features,
the predictions, and the two CRSP-derived crosswalks whose withdrawal
Appendix~\ref{app:ethics} records. Nothing is held back from the assessors in
particular: no reader receives those rows, and any holder of a CRSP licence
regenerates them exactly from the released scripts.
Sections~\ref{app:external:benchmarks} to~\ref{app:external:neural} then record
external material of a different kind: where the benchmark built here stands
among the corpora the background survey's literature was built on, and one
branch of price forecasting that survey passes over.

\FloatBarrier
\section{Released Artefact Inventory}
\label{app:external:release}

The repository's public core is the directory \texttt{release/}, engineered so that a
third party can reconstruct the exact study sample without receiving a single
licensed market-data value. Two CRSP-derived files an earlier staging
carried are withheld from it, as Appendix~\ref{app:ethics} sets out. Auditing the reported tests is the wider
repository's job, not this directory's: the current evidence tables live under
\texttt{results/tables/}, and the \texttt{aggregate\_results/} bundle inside
\texttt{release/} is the pre-day-clustering vintage the earlier analysis quoted,
as its own row below states. The
membership and link crosswalks an earlier staging shipped carry CRSP identifiers and dates, on the
reading Appendix~\ref{app:ethics} records.
Table~\ref{tab:app-release} is its index.

\begin{table}[htbp]
  \centering
  \caption[The released artefact inventory]{The released artefact inventory (\texttt{release/} in the repository).}
  \label{tab:app-release}
  \small
  \begin{tabular}{lp{9.4cm}}
    \toprule
    Artefact & Contents \\
    \midrule
    \texttt{accession\_index.csv} & Accession index of all 144{,}129 benchmark filings: accession number, CIK, ticker, form, filing timestamp, effective trading day and split. Doubles as the exact SEC~EDGAR pull list, so no filing text needs redistributing \\ % src: release/accession_index.csv (144,129 data rows, verified 2026-08-02); release/DATA_CARD.md
    \texttt{membership\_intervals.csv} & The 914 survivorship-free point-in-time S\&P~500 membership intervals (ticker, PERMNO, CIK, entry and exit dates). \emph{Withheld from the public core as CRSP-derived; rebuilt by the ingestion script} \\ % src: crsp_restricted/membership_intervals.csv (914 data rows, verified 2026-08-02)
    \texttt{cik\_links\_pit.csv} & The 30{,}100 point-in-time PERMNO$\to$CIK link windows, including the 31 audited successor-CIK override windows covering 23 firms that recover filings from acquired and renamed registrants. \emph{Withheld on the same ground} \\ % src: crsp_restricted/cik_links_pit.csv (30,100 data rows, verified 2026-08-02); release/DATA_CARD.md
    \texttt{split\_definition.md} & The chronological split rule and its filing counts (long-form 19{,}668/3{,}963/7{,}970 and event-driven 73{,}088/14{,}266/25{,}174 across train/validation/test), re-derivable from the effective trading day alone \\ % src: release/split_definition.md
    \texttt{DATA\_CARD.md} & The datasheet: motivation, composition, collection, recommended uses, exclusions and the required CRSP/WRDS attribution \\ % src: release/DATA_CARD.md
    \texttt{aggregate\_results/} & A twelve-file bundle of observation-level standalone comparisons, stratified breakdowns and the gate-cushioning table, cut before the inference stack moved to day clustering and seed ensembling; it is retained as the vintage the earlier analysis quoted. The tables this dissertation cites are the working evidence set under \texttt{results/tables/}, which travels with the repository rather than inside \texttt{release/}. Both are summary statistics only, carrying no per-row licensed value \\ % src: release/README.md; release/aggregate_results/
    \texttt{config\_hashes.csv}, \texttt{run\_configs/} & The fingerprint chain: one row per production run and the 240 \texttt{config.json} preimages themselves, so the fingerprint check is executable rather than asserted \\ % src: release/config_hashes.csv (240 data rows, verified 2026-08-02); release/run_configs/ (240 files)
    \texttt{tag\_manifest.txt} & The 26 repository tags, each with the commit it pins, that commit's date and the date the tag was written: the 21 pre-registration tags in eight families, one pre-test protocol and four writing milestones \\ % src: release/tag_manifest.txt (26 tag rows)
    \texttt{raw\_generations/} & 1{,}221 parquet shards holding 608{,}221 prompted-model generations (each row the model's verbatim output string beside its parsed horizon forecasts), with the prompt templates and decoding configuration that produced them \\ % src: release/raw_generations/README.md
    \bottomrule
  \end{tabular}
\end{table}

Four further components ship alongside that index. The \emph{licence-free
variant} rebuilds the price side from public sources so that the benchmark can be
exercised end to end with no subscription at all. Its coverage is priced honestly
rather than advertised; Section~\ref{sec:data-public} states what it costs.
The \emph{fingerprint chain} covers 240 production runs carrying 233 distinct
SHA-256 configuration digests, the 7 collisions being expected pairs whose
configurations are byte-identical. % src: results/tables/config_fingerprints.md
Its index publishes two digests per run: the original fingerprint and the
fingerprint of the shipped preimage. Both are published because 38 preimages had
absolute local paths neutralised before release, each row naming the exact
fields that moved, while the remaining 202 hash identically to their originals. The released script
\texttt{config\_fingerprints.py -{}-verify-preimages} re-hashes every preimage
against the index using only files inside the repository. % src: release/config_hashes.csv (240 rows; 38 with a non-blank sanitised_fields, 202 whose released_config_sha256 equals config_sha256)
The \emph{test suite} ships in full: 186 tests in 33 test modules, all passing, covering
ingestion, alignment, split assignment, training and evaluation, and including
byte-identity checks that the pre-tokenisation text caches reproduce their
uncached reference exactly. A reader counting \texttt{def test\_} by hand will
find 184: two of those definitions are parametrised over two model classes each,
so the collected count is 186 and the file count is
exact. % src: writing/dissertation/supporting/REPRODUCIBILITY.md; tests/ (33 modules, 184 test functions, two @pytest.mark.parametrize decorators in tests/models/test_classical_text.py, no conftest.py and no class-based tests)
Finally the seven \emph{registered analysis records} (\texttt{configs/prereg\_*.md})
travel with the code, so a reader can check each audit family's pre-declared
design against what was reported. % src: configs/prereg_*.md (7 files, verified 2026-08-02); release/README.md

% The inventory table is flushed here so that it cannot be deferred into the
% directory listing below, which it interrupted mid-tree in earlier builds.
\clearpage

\FloatBarrier
\section{Repository Layout}
\label{app:external:layout}

The code is one importable Python package, \texttt{sp500vol}, under
\texttt{src/}, driven by thin command-line scripts and declarative
configurations:

\begin{verbatim}
.                       the repository root
|-- configs/            declarative experiment specifications
|   |-- base.yaml         shared defaults: paths, seeds, split boundaries
|   |-- data/             dataset scopes (full, sample, dry-run)
|   |-- models/           one YAML per configured arm (A1-A5, B1-B4, C1-C5,
|   |                       D1-D3); A6, A7, C6 and D4 are script-driven
|   |-- ablations/, compute/   sweep and hardware profiles
|   `-- prereg_*.md       the seven registered analysis records
|-- data/universe/      scaffolding; the CRSP crosswalks are withheld
|-- release/            the public bundle inventoried above
|-- results/
|   |-- runs/<run_id>/    per-run artefacts; the predictions are withheld
|   `-- tables/           the evidence tables named by this report's sources
|-- scripts/            command-line entry points
|   |-- ingest_wrds.py    CRSP export -> membership, PIT links, returns store
|   |-- build_dataset.py  EDGAR fetch and parse -> aligned filing-horizon rows
|   |-- train.py, evaluate.py, compare_runs.py, dm_vs_baseline.py
|   |-- analysis/         the audit layer: clustered_dm.py,
|   |                     forecast_combination.py, firm_identity_control.py,
|   |                     anon_*.py, matched_firm_swap.py,
|   |                     signal_injection_power.py, config_fingerprints.py,
|   |                     supp_figs/ (gated figure generators)
|   |-- experiments/      arm drivers: prompted-LLM (incl. prompt.py),
|   |                     tuning arm, external-panel replications
|   `-- release/          bundle staging, disclosure scan, package build
|-- src/sp500vol/       the package
|   |-- data/             universe, crsp, edgar_fetcher, parser,
|   |                     trading_calendar, alignment, pipeline_state
|   |-- features/         returns, volatility, text_preprocess, lm_dictionary
|   |-- models/           price/, classical_text/, neural_text/, fusion/
|   |-- evaluation/       metrics, dm_test, bootstrap
|   |-- training/         end-of-epoch checkpointing for resume on preemption
|   `-- utils/            seeding, cost tracking, environment capture
|-- tests/              the 186-test suite
`-- writing/            this report's LaTeX sources and frozen drafting sources
\end{verbatim}

The division of labour is deliberate. \texttt{src/} holds everything a model or a
metric needs and nothing about a particular experiment; \texttt{configs/} holds
every experimental decision as data, which is what makes the fingerprint chain
meaningful. \texttt{scripts/} holds the executable pipeline, with
\texttt{scripts/analysis/} as the audit layer that produced the reference ladder,
the identity controls, the power calibration and the figures. Every artefact
under \texttt{results/tables/} that this report cites is written by a committed
generator, never by hand, and each generator aborts if its source table has
drifted from the values it recorded. Four are not, and are named here rather
than left to be discovered: \texttt{event\_driven\_text\_cache\_validation.txt},
\texttt{gate\_cushioning.md}, \texttt{rtx3090\_probe\_24g\_summary.md} and
\texttt{yelp\_cascade\_synthetic.md} are records of one-off validation or pilot
runs whose driver was never committed, so they can be read but not
regenerated. % src: a search of scripts/, src/ and the Makefile for each basename under results/tables/; these four return nothing, while the dm_* and compare_* families are written under names their generators build at run time

\FloatBarrier
\section{Regenerating the Study from Scratch}
\label{app:external:regen}

Two external inputs must be supplied by the regenerator, because neither can be
redistributed with the repository. The first is a \textbf{CRSP/WRDS export}
\cite{crsp2026data}: survivorship-free S\&P~500 index-constituent membership,
daily split- and dividend-adjusted total returns for 2010--2025, and the
CRSP/Compustat-Merged link table used for point-in-time PERMNO-to-CIK resolution.
The second is the \textbf{Loughran--McDonald master dictionary}
\cite{loughran2011liability}, required by the two dictionary arms and obtained
from its authors' distribution site. Filing text itself needs no supply: it is
fetched from SEC~EDGAR by accession under the published fair-access limits
\cite{sec2024edgar}, and the released accession index is the exact pull list.
Pretrained checkpoints \cite{araci2019finbert,beltagy2020longformer,devlin2019bert,liu2019roberta,qwen3_2025}
are downloaded through the usual model hub \cite{wolf2020transformers}; the
software stack is PyTorch \cite{paszke2019pytorch}, scikit-learn
\cite{pedregosa2011sklearn} and, for the prompted arm, offline vLLM
\cite{kwon2023vllm}.

Regeneration then runs in four stages: \texttt{ingest\_wrds.py} builds the
membership table, the point-in-time link table and the returns store;
\texttt{build\_dataset.py} fetches and parses EDGAR text and emits the aligned
filing-by-horizon rows under the no-look-ahead rule of
Section~\ref{sec:data-labels}; \texttt{train.py} trains one arm on one disclosure
subset for one seed; and \texttt{evaluate.py}, \texttt{compare\_runs.py} and
\texttt{dm\_vs\_baseline.py}, followed by the \texttt{scripts/analysis/} layer,
rebuild the metric, inference and audit tables.
Figure~\ref{fig:app-pipeline} draws the pipeline those stages rebuild, from the
three sources through to the reproducibility record, and
Figure~\ref{fig:app-components} draws the same programme as code: which
packages the library holds, how large each is, and the only four import edges
between them.
Appendix~\ref{app:hyperparams}
gives the exact per-arm command templates and settings. Budgeting matters here:
the neural training fleet's 144 production runs consumed 590.4 GPU-hours on
NVIDIA A100-SXM4-40GB hardware. % src: results/tables/cost_accuracy.md; TECHNICAL_SUPPLEMENT.md compute scope note; results/tables/a100_40g_probe_summary.md
A full replication is therefore a fleet-scale undertaking. Re-deriving the audit
layer on top of it is CPU-only and cheap, but it cannot be done from the
released aggregate tables: those are the layer's outputs, and its inputs are the
per-observation loss differentials in the per-run prediction files that
Section~\ref{app:external:withheld} withholds. One
environment freeze survives with declared scope: \texttt{ENVIRONMENT\_inference.txt},
224 packages, pinning the prompted-inference stack rather than the
encoder-training stack, whose declared constraints live in
\texttt{pyproject.toml}. % src: writing/dissertation/supporting/REPRODUCIBILITY.md

% The schematic is drawn inline rather than imported so that it cannot drift
% from the sections it depicts. Every box states a stage sec:data-sources or
% sec:design already specifies; no stage is added here.
\begin{figure}[htbp]
  \centering
  \begin{tikzpicture}[
    font=\footnotesize,
    box/.style={draw, rounded corners=1.5pt, align=left, text width=66mm,
                inner xsep=2.2mm, inner ysep=1.8mm},
    head/.style={align=center, text width=70.4mm, font=\footnotesize\bfseries},
    flow/.style={draw, -{Latex[length=1.8mm,width=1.4mm]}, semithick}]

  \node[head, anchor=north] (h1) at (0mm,0mm)
    {The benchmark (Section~\ref{sec:data-sources})};
  \node[box, below=1.6mm of h1] (a1)
    {\textbf{Sources.} CRSP index constituents and daily total returns, and the
     CRSP/Compustat-Merged link table, licensed through WRDS; the SEC EDGAR
     full-text archive.};
  \node[box, below=3mm of a1] (a2)
    {\textbf{Point-in-time universe.} 914 S\&P~500 membership intervals over
     2010--2025 on 832 PERMNOs, delisted members retained rather than
     dropped.}; % src: crsp_restricted/membership_intervals.csv (914 intervals, 832 PERMNOs)
  \node[box, below=3mm of a2] (a3)
    {\textbf{Entity linking.} PERMNO to CIK resolved window by window through
     the CCM link table: 30{,}100 point-in-time windows, including 31 audited
     successor overrides covering 23 firms.}; % src: crsp_restricted/cik_links_pit.csv; release/DATA_CARD.md lines 45-46
  \node[box, below=3mm of a3] (a4)
    {\textbf{Disclosure corpus.} 10-K, 10-Q and 8-K text fetched from EDGAR
     under the resolved CIK: 144{,}443 filings parsed, 144{,}129 aligned to the
     benchmark (31{,}601 long-form, 112{,}528 event-driven).}; % src: frozen 03_dataset.tex; TECHNICAL_SUPPLEMENT.md Table S1
  \node[box, below=3mm of a4] (a5)
    {\textbf{Labels and price features.} Forward realised volatility at
     $h=5$, 10 and 20 trading days; every filing mapped to its effective
     trading day; HAR components ending at the last close before the filing;
     no-look-ahead audit reporting zero violations.}; % src: frozen 03_dataset.tex (horizons; effective trading day; zero violations)
  \node[box, below=3mm of a5] (a6)
    {\textbf{Modelled panel and splits.} 431{,}245 aligned filing-by-horizon
     rows, 427{,}429 surviving the feature--label joins; chronological splits,
     training 2010--2019, validation 2020--2021, test 2022--2025.}; % src: TECHNICAL_SUPPLEMENT.md Table S1; release/split_definition.md

  \node[head, anchor=north] (h2) at (78.4mm,0mm)
    {The audit (Section~\ref{sec:design})};
  \node[box, below=1.6mm of h2] (b1)
    {\textbf{Model matrix.} Block~A price references, Block~B lexical
     baselines, Block~C contextual encoders and large language models,
     Block~D price--text fusion.};
  \node[box, below=3mm of b1] (b2)
    {\textbf{Training.} One fixed recipe for every trained Block~C and~D arm,
     three seeds (2026--2028), and the per-observation seed ensemble as the
     primary forecast.}; % src: frozen 05_protocol.tex; seed_aggregate.md
  \node[box, below=3mm of b2] (b3)
    {\textbf{Forecast combination.} Price and text nested in one log-space
     combiner, every weight fitted on the validation years and frozen before
     the test years are touched.};
  \node[box, below=3mm of b3] (b4)
    {\textbf{Reference ladder.} The recalibrated HAR, then a
     firm-identity-augmented reference, then a maximal price pool, then the
     conjunction of all three.};
  \node[box, below=3mm of b4] (b5)
    {\textbf{Inference.} Day-clustered Diebold--Mariano tests, Holm correction
     inside 15 pre-declared families, a label-shuffle placebo, and a
     signal-injection power calibration.}; % src: results/tables/holm_families.md (15 families); results/tables/signal_injection_power.md
  \node[box, below=3mm of b5] (b6)
    {\textbf{Economic adjudication.} A portfolio Sharpe test, a value-at-risk
     backtest and a utility-based fee, each fitted on validation and frozen on
     test.};
  \node[box, below=3mm of b6] (b7)
    {\textbf{Reproducibility.} Self-describing run directories, 240
     fingerprinted production runs and a 186-test suite, with generators that
     abort when a source table drifts.}; % src: results/tables/config_fingerprints.md (240 runs); writing/dissertation/supporting/REPRODUCIBILITY.md (186 tests)

  \foreach \i/\j in {a1/a2, a2/a3, a3/a4, a4/a5, a5/a6,
                     b1/b2, b2/b3, b3/b4, b4/b5, b5/b6, b6/b7}
    \draw[flow] (\i.south) -- (\j.north);
  \draw[flow] (a6.east) -- ++(4mm,0mm) |- (b1.west);
  \end{tikzpicture}
  \caption[The pipeline behind SP500Vol-Bench and its audit]{The pipeline of
  Chapter~\ref{ch:data}, drawn as a flow. The left column is the benchmark of
  Section~\ref{sec:data-sources}, running from the three sources to the
  modelled panel and its chronological splits. The right column is the audit of
  Section~\ref{sec:design}, running from the model matrix to the
  reproducibility record. Every box states a stage one of those two sections specifies,
  and every count is the one that section reports; nothing is drawn that they
  do not state. The two external portability panels of
  Section~\ref{sec:data-portability} sit outside this pipeline and are not
  drawn. The diagram carries no result. It shows what was run and in what
  order, not what was found.}
  \label{fig:app-pipeline}
\end{figure}


\begin{figure}[tbp]
  \centering
  \begin{tikzpicture}[
    font=\footnotesize,
    pkg/.style={draw, rounded corners=1.5pt, align=left, text width=42mm,
                inner xsep=2.0mm, inner ysep=1.6mm},
    grp/.style={draw, dashed, rounded corners=2pt, inner xsep=2.6mm, inner ysep=2.6mm},
    dep/.style={draw, -{Latex[length=1.8mm,width=1.4mm]}, semithick},
    lab/.style={align=center, font=\footnotesize\bfseries}]

  \node[lab] (h) at (0mm,5mm) {The library: \texttt{src/sp500vol} --- 60 modules, 8{,}864 lines};

  \node[pkg, anchor=north west] (data) at (-46mm,0mm)
    {\textbf{data} \hfill 9 files, 1{,}574 lines\\
     EDGAR fetching, parsing, point-in-time alignment, split assignment.};
  \node[pkg, anchor=north west] (models) at (4mm,0mm)
    {\textbf{models} \hfill 32 files, 6{,}366 lines\\
     every arm behind one base class: \texttt{price} (6),
     \texttt{classical\_text} (7), \texttt{neural\_text} (13),
     \texttt{fusion} (4).};

  \node[pkg, anchor=north west] (feat) at (-46mm,-30mm)
    {\textbf{features} \hfill 5 files, 239 lines\\
     log returns, the annualised label, text preparation.};
  \node[pkg, anchor=north west] (utils) at (4mm,-30mm)
    {\textbf{utils} \hfill 6 files, 327 lines\\
     seeding, environment capture, run paths, hashing.};

  \node[pkg, anchor=north west] (eval) at (-46mm,-52mm)
    {\textbf{evaluation} \hfill 4 files, 217 lines\\
     losses and the Diebold--Mariano machinery. Imported by nothing in the
     library; the scripts call it.};
  \node[pkg, anchor=north west] (train) at (4mm,-52mm)
    {\textbf{training} \hfill 2 files, 138 lines\\
     checkpoint save and resume. \textbf{pipelines} is declared and empty.};

  \draw[dep] ([xshift=-10mm]data.south) -- ([xshift=-10mm]feat.north);
  \draw[dep] ([xshift=10mm]models.south) -- ([xshift=10mm]utils.north);
  \draw[dep] ([xshift=8mm]data.south) -- ([xshift=-12mm]utils.north);
  \draw[dep] ([xshift=-8mm]models.south) -- ([xshift=12mm]feat.north);
  \end{tikzpicture}
  \caption[The library, component by component]{The same programme as
  Figure~\ref{fig:app-pipeline}, read as code rather than as a flow. Counts are
  tracked \texttt{.py} files and their lines. Arrows are the only four
  package-level import edges the source contains, derived by parsing every
  import statement rather than inferred from the names. Three non-edges are
  worth as much as the four edges: \texttt{models} imports neither
  \texttt{data} nor \texttt{evaluation}, so an arm cannot reach the corpus or
  its own score while it runs, and nothing under \texttt{src/} imports
  \texttt{scripts/}. Inside \texttt{models} the four arm families each depend
  only on the shared base class, except that \texttt{fusion} and
  \texttt{neural\_text} reuse the text loader in \texttt{classical\_text}. The
  test suite is a sibling of the package, not a part of it: 33 modules
  collecting 186 cases, nine of which exercise \texttt{scripts/} rather than
  the library.} % src: parsed from src/sp500vol/ (60 tracked .py files; 4 package-level import edges: data->features, data->utils, models->features, models->utils); tests/ (33 modules, 184 test functions, two parametrised over two model classes each, so 186 collected)
  \label{fig:app-components}
\end{figure}

\FloatBarrier
\section{Withheld Data, and Why Withholding Is Reproducibility-Neutral}
\label{app:external:withheld}

The CRSP/WRDS licence permits academic research use and bars redistribution of
the raw data and of per-row quantities derived from it \cite{crsp2026data}.
Withheld, therefore, are the daily return series, the realised-volatility labels,
the price features and the per-row model predictions: the last because a
predictions file embeds its label in plain form and is licensed data by
construction. Per-row loss-differential series could in principle be released in
a non-invertible form, carrying differentials alone with no predictions, labels
or features from which a label could be reconstructed. Release would require
written confirmation from the institutional licensing administrator, which the
project did not seek to assume.

The withholding is nonetheless reproducibility-neutral, by design rather than by
assertion. Everything needed to \emph{regenerate} the withheld rows is
released: the universe, the point-in-time links, the split rule, the labelling
code and the pinned configurations. What comes back is not uniform across that
set, and the difference is stated rather than glossed. The CRSP-derived rows
return exactly: the label-parity gate recomputes all 431{,}245 aligned
realised-volatility labels from the raw series on the committed windows and
finds them bitwise identical. The per-row model predictions do not, and this
report turns on that fact elsewhere --- the two executed anonymisation arms were
retained through the registered deviation route precisely because their
regenerations were not bit-identical, one matching 111{,}458 of 117{,}407
predictions and the other none of 333{,}192. A subscription holder therefore
rebuilds the licensed inputs byte-for-byte and the GPU-produced outputs only to
within the nondeterminism this report documents.
% src: results/tables/label_parity.md GATE 1 (n = 431,245, unreconstructed 0, max abs diff 0.000e+00, bitwise exact True); results/tables/anon_arm.md (G1 c6 exact 111,458/117,407; G1 c2 exact 0/333,192)
Everything needed to \emph{audit} the reported claims without a
subscription is released too, because the aggregate evidence tables are summary
statistics carrying no per-row licensed value. Every test statistic, p-value and
correction printed in Chapters~\ref{ch:results} and~\ref{ch:validation} appears
in them. And the one arm whose inputs are entirely public is the prompted model,
which consumes SEC text and emits JSON. It ships nearly complete: prompts, decoding
configuration and 608{,}221 generations, so those outputs can be re-parsed and
re-scored by anyone. The exception is named in the repository rather than left to
be found: the date-only and date-plus-ticker probe generations were not
retained, and their results survive only in the aggregate evidence tables. % src: release/raw_generations/README.md ("The Qwen date-only and date+ticker probe generations were not retained; their results survive in the released aggregate evidence tables.")

\FloatBarrier
\section{Verifying the Numbers in This Report}
\label{app:external:verify}

Three mechanisms make this report checkable. First, the \textbf{source
comments}: every number in the \LaTeX{} sources carries a same-line
\texttt{\% src:} comment naming where it was read from. Of the 868 such comments, 799 name a file the repository carries, so a
marker opening the sources goes from the printed number to the artefact behind
it in one step. Forty-two of those lead with a file in \texttt{release/}, the
public core; the rest lead with results tables, configurations, scripts and
frozen drafting sources. A further 54 name an artefact the licence withholds ---
the per-run predictions and the two CRSP-derived crosswalks --- and one is a
wildcard standing for four configuration files that do ship. The
remaining 14 the census cannot resolve to one file: ten name a cited paper, three
the repository's own commit history and one this report's own chapters. The census is regenerated
from the sources rather than typed, and the build fails if these counts drift.
% src: results/tables/src_comment_census.txt (scripts/analysis/src_comment_census.py)
Every count this report leads
with --- the standalone census, each rung of the ladder, the injection
recoveries and the economic adjudication --- names a file in
\texttt{results/tables/} at its point of use. Any number carried on a secondary basis (a single seed, an
observation-level test, a variance-unit restatement) says so in the prose
beside it. Second, the \textbf{run-artefact contract}: every production run
writes a self-describing directory named
\texttt{\{model\}\_\{dataset\}\_\{disclosure\}\_seed\{seed\}} containing the
fully resolved configuration (\texttt{config.json}), row-level predictions and
per-split metrics. All 240 fingerprinted runs carry those three. An environment
snapshot (\texttt{env.json}) and wall-clock and GPU-hour accounting
(\texttt{cost.json}) are carried by 183 of the 240. The 57 without them are the
38 prompted runs, which train nothing; the six SHAR and HARQ price runs; the six
anonymisation control and mask legs, the four TF--IDF ones being the lineage
whose missing package versions Section~\ref{sec:val-repro} records; and seven
ancillary probe arms. The cut does not fall on training. Of the 178
training-bearing runs 171 carry all five, the seven that do not being those
anonymisation legs and the delta-text probe; and twelve runs that train nothing
at all --- the A1, A3, A4 and A5 price models across the three channels --- carry
them.
% src: release/config_hashes.csv against results/runs/*/ (240 with config.json,
% metrics.json and predictions.parquet; 183 with env.json and cost.json;
% training_sha256_or_blank non-blank in 178, 18 distinct, 171 of them with env.json)
Appendix~\ref{app:fullresults}
reads its grids mechanically from stored artefacts rather than transcribing them
from logs. Third, the \textbf{drift gates}: the fingerprint script re-hashes the
released configurations against the published index, and every table and figure
generator aborts when its source table no longer matches the numbers it
recorded. A stale artefact therefore cannot quietly survive an upstream change.
The recommended order for a marker with the repository to hand is to run
\texttt{config\_fingerprints.py -{}-verify-preimages}, run the test suite, then
pick any claim in Chapter~\ref{ch:results}, read its \texttt{\% src} comment, and
open the named table in \texttt{results/tables/}, where 113 of Chapter~\ref{ch:results}'s 150 source comments lead.

One reconciliation those three mechanisms do not perform for the reader, and
which a careful marker will notice, is worth stating here rather than leaving
to be found. Panel~(c), Figure~\ref{fig:data-membership-b}, stratifies the
long-form test panel by whether the filer later exits, on 822 exiting and
7{,}106 survivor filings. Those sum to 7{,}928, which is 42 fewer than the
7{,}970 long-form test filings of Table~\ref{tab:data-splits}. The two are not
the same denominator: that panel is carried on the superseded early-campaign
convention, as its own caption states, so its row support is that campaign's
rather than the released split's. Nothing else in
the report is computed on the 7{,}928, and no count anywhere is taken as a
share of it.

\FloatBarrier
\section{Project Management, Version Control and Continuous Integration}
\label{app:external:process}

The preceding sections describe what the project produced. This one describes
how it was run, because a reproducibility claim is only as strong as the process
that generated the artefacts behind it, and that process is itself inspectable.
Two things are kept apart throughout. What the process \emph{was} is a matter of
record in the tag list, the workflow definition and the build files; what it
could and could not guarantee is an argument, and it is made explicitly so that
a reader can discount it. The tag list is a property of the version-control
history rather than of any single file, so it is enumerated below in full.

\textbf{Planning: evidence-gated increments, not time-boxed sprints.} The
project had one author and no external client, so no stakeholder existed from
whom requirements had to be elicited. The requirements are the research question
of Section~\ref{sec:intro-aim}, decomposed into the five objectives of
Section~\ref{sec:intro-objectives}, and every deliverable in
Section~\ref{sec:intro-deliverables} answers to one of them. No sprint cadence,
burndown chart or ticket board was kept, and none is reconstructed here: a
borrowed vocabulary applied after the fact would describe this project less
accurately than what actually governed it.

\textbf{Risk: what was not registered, and what actually materialised.} No risk
register was kept either, and none is reconstructed here for the same reason.
What can honestly be given is the other half: the risks that did materialise,
the control that was standing against each when it did, and what the control
cost. Every one is documented elsewhere in this report; gathering them in one
place is what a register would have been for.

\emph{Licence exposure.} The market data are licensed for research and may not
be redistributed, which threatened the release rather than the result. The
control was designed in rather than added: per-row licensed quantities are
withheld and fully regenerable by any licence holder from the released scripts,
and a licence-free variant rebuilds prices from public sources at 72.6, 89.3 and
95.2 per cent coverage of the three splits. The cost is that the public core is
not the panel the results are computed on, which
Section~\ref{sec:data-sources} prices rather than hides.

\emph{Lawful access.} A bulk retrieval from a free public service is a standing
risk to the service and to the project's access. The fetcher declares its
identity, stays inside the published rate limits, and logs every request in a
resumable, retry-safe ledger so that no document is requested twice. The control
was in place before the first bulk pull, not after a failure.

\emph{Silent survivorship bias.} This one materialised. Where a constituent was
acquired or reincorporated, the vendor link could point a historical identifier
at the successor's registrant, under which the original firm filed nothing
during its membership window --- a failure that removes exactly the delisted and
acquired firms whose retention the benchmark's central claim rests on, and
removes them without an error. It was caught by audit rather than by inspection,
and is now held by 31 audited override windows covering 23 firms, a regression
guardrail and dedicated unit tests; the no-look-ahead audit reports zero
violations for those 23 firms in
particular. % src: crsp_restricted/cik_links_pit.csv (30,100 link windows, 31 override windows over 23 firms); frozen 03_dataset.tex (no-look-ahead audit, zero violations including the 23 override firms)

\emph{Loss of the compute estate.} A fleet-scale campaign on rented hardware
carries the risk that the environment disappears before the analysis is
finished, and this one materialised twice. The FinBERT checkpoints were
destroyed in a disk cleanup, which cost the long-form matched-swap replication
one of its three arms, the other two having missed their own reproduction gates
for unrelated reasons; and the earliest run metadata logged no package versions,
which retired three further arms --- the long-form anonymisation exit, the
event-driven anonymisation and the swap TF--IDF arm --- at their gates rather
than allowing them to be rebuilt. The control that was standing --- the reproduction gate
itself, and a written retirement record for every arm that failed one --- did
not save the arms, and was never going to. What it did was make the loss
visible: the retirements are on the record with their thresholds and measured
deviations, and Section~\ref{sec:conc-limitations} states what the missing arms
would have measured. An artefact-retention policy and an environment snapshot on
every run from the first day would have prevented both, and their absence is the
one shortfall in this project's engineering that is hygiene rather than
scope. % src: results/tables/swap_longform_retirement.md (gate, threshold and measured deviation per retired arm); results/runs/*/env.json (183 of 240 carry a snapshot)

The pattern is worth stating because it is not flattering throughout. Two of the
four were designed against in advance and held. One was caught by an audit that
existed for a different purpose. One was not anticipated at all, and the only
control that applied was the one that made the damage countable.

What governed it was the \emph{arm}: a named experimental or audit unit that was
specified in writing, registered under a version-control tag before it ran,
executed, and then either admitted to the evidence base or retired at a stated
gate. Increments closed on evidence rather than on the calendar, and the gates
that closed them are the ones the preceding chapters already document. A derived
arm's statistic may be cited only after a rerun reproduces that arm's committed
predictions to a relative tolerance of $10^{-8}$. A cross-family probe counts as
a live instrument only if it passes the forecaster-health screen, whose two
thresholds read the run's own forecasts and never the contrast under test
(Section~\ref{sec:val-instruments}). And a combination cell is declared genuine
only under the conjunction of clustered significance, Holm correction and the
label-shuffle placebo (Section~\ref{sec:design-inference}). The honest record of
a scheme like this is not its list of successes but its register of failures,
and Section~\ref{sec:val-repro} reports that register in full: the long-form
matched-swap replication lost all three of its arms, and the long-form
anonymisation channel closed with none executed, each exit carrying its gate,
its threshold and its measured deviation. One further registered increment was
retired without ever being scored: the fusion gate's internal weight readout,
whose evidence in this report is behavioural because the direct readout needs a
forward pass that was never run once the fleet was decommissioned. It is one
pass for anyone holding the checkpoints, and it is recorded here rather than
among the open questions of Section~\ref{sec:conc-future} because what the gate
suppressed moves no count. % src: TECHNICAL_SUPPLEMENT.md S19 register
An increment-based plan whose abandoned increments went unreported would be
indistinguishable from no plan at all.

\textbf{What the version-control record shows, and what it cannot.} The
repository carries 26 tags. Twenty-one of them mark pre-registrations in eight
families: seven tags for the long-form matched swap and the entity-anonymisation
arm (\texttt{prereg-ea-v1.0} to \texttt{v1.6}), four for the residual family
audit (\texttt{prereg-rfa-v1.0} to \texttt{v1.3}), three for the earnings-call
panel audit (\texttt{prereg-maec-v1.0} to \texttt{v1.2}), two for the
elicitation-fairness test (\texttt{prereg-ef-v1.0} and \texttt{v1.1}), two for
the tuned challenger arm (\texttt{hpo-prereg-v1.0} and \texttt{v1.1}), and one
each for the mechanism-and-relabelling record (\texttt{prereg-cd-v1.0}), the
licence-free public-price variant (\texttt{prereg-h-v1.0}) and the founding 10-K
corpus audit (\texttt{prereg-kc-v1.0}). % src: git tag -l (26 tags; 21 pre-registration tags in exactly the eight families named, verified 2026-08-04)
The remaining five mark one further pre-test protocol and four writing
milestones. The seven \texttt{prereg-} families correspond one to one with the
seven registered analysis records shipped under \texttt{configs/}. That mapping is
checkable rather than asserted without a copy of the version-control history:
\texttt{release/tag\_manifest.txt} lists every tag beside the commit it pins and
its date, so the eight families can be counted from the repository alone. A second
version number inside a family records an amendment carried by its own tag rather
than a silent edit to the registered file. % src: configs/prereg_*.md (7 records: swap_lf_and_anon, residual_family_audit, maec_audit, elicitation_fairness, mechanism_and_labels, public_variant, kogan_corpus)

A tag is strong evidence of one thing and no evidence of another, and conflating
the two would be an overclaim of exactly the kind this report exists to resist.
It fixes the bytes of the specification the author later claims to have
followed, so any divergence between what was registered and what was reported
can be found mechanically by checking the tag out and comparing. It does not
establish temporal priority. The author holds the repository and the clock
alike, a tag can be moved and a history can be rewritten, so an internal tag
evidences internal consistency and not that the specification preceded the
result. That limitation is why one arm, and only one, carries the external
instrument described at the end of this section.

\textbf{Continuous integration: what runs, and what is only advised.} A workflow
at \texttt{.github/workflows/ci.yml} runs on every push to \texttt{main} and on
every pull request opened against it. Its first job installs the locked
dependency set and runs the test suite on Python 3.11 under the marker
expression \texttt{not slow and not gpu and not network}, so the fast
correctness tests execute on every push while the fleet-scale, GPU-bound and
network-dependent tests are excluded by design. Its second job holds the library and its tests to the
full lint rule set, running \texttt{ruff check} and \texttt{ruff
format -{}-check} over \texttt{src/} and \texttt{tests/}; it checks
\texttt{scripts/} with \texttt{ruff check -{}-select F821,F811,E9}, which is
the errors that would make an analysis script wrong --- an undefined name, a
redefinition, a syntax error --- and not its style, because restyling code
whose outputs are frozen would rewrite files this report depends on for no
gain. It then runs \texttt{mypy} over
\texttt{src/sp500vol/}. % src: .github/workflows/ci.yml (triggers push and pull_request on main; python-version matrix ["3.11"]; pytest marker expression; two ruff lint scopes, one ruff format scope, mypy step)

Two properties of that arrangement need stating plainly, because a flattering
description of them would be the very error this report audits others for.
First, the type-checking step is invoked as
\texttt{mypy src/sp500vol/ \textbar\textbar\ true}, beside a comment recording
that strict typing was deferred, so static type checking here is \emph{advisory
and not enforced}: a type error cannot fail the run. The \texttt{lint} target of
the \texttt{Makefile} runs the same checker without that escape, so the check is
available on demand, but continuous integration does not gate on
it. % src: .github/workflows/ci.yml (mypy step trailing `|| true`, with the in-file comment deferring strict typing); Makefile (lint target runs mypy with no trailing escape)
Second, the project has a single author and its history contains no merge
commits at all, so although the workflow declares a pull-request trigger, no
review by a second party ever took place. What the workflow supplies is a
regression check on one person's pushes, and describing it as a peer-review
process would be false. % src: git log --format='%an' (one author across the whole history); git log --merges (0 merge commits, verified 2026-08-04)

\textbf{Dependency and environment reproducibility.} Three mechanisms operate at
three time scales. \texttt{pyproject.toml} declares the constraints, including a
floor of Python 3.11; \texttt{uv.lock} pins the fully resolved dependency graph
those constraints admit, and the workflow installs through \texttt{uv sync}
against that lock file rather than re-resolving, so the environment that is
checked is the environment that is declared. The \texttt{Makefile} exposes the
routine operations as named targets, among them \texttt{make test},
\texttt{make lint}, \texttt{make format}, \texttt{make data},
\texttt{make train}, \texttt{make tables} and \texttt{make dissertation}, so the
command that produced an artefact is committed rather than
remembered. % src: pyproject.toml (requires-python >= 3.11); uv.lock; Makefile (target list)
The third mechanism is the per-run environment snapshot of
Section~\ref{app:external:verify}, carried by 183 of the 240 runs on the coverage
that section sets out, with the further exception Section~\ref{app:external:regen}
states. The retirement register of
Section~\ref{sec:val-repro} is what that third mechanism looks like when it
fails: the arms whose early snapshots recorded no package versions could not be
rebuilt, and were retired rather than reconstructed by guesswork.

\textbf{The external timestamp, scoped to the one arm that needed it.} One arm
is exposed enough to the priority objection above to warrant an
author-independent instrument.
The tuned challenger of Section~\ref{app:hyperparams:asha} exists to
test whether the fixed training recipe rather than the text produced the null,
and a hyperparameter search is the one arm whose specification, had it been
free to follow the test evaluation, could have been steered towards whichever
conclusion its author preferred. Its pre-registration bundle was therefore
stamped outside the repository.

The stamped object is a single archive,
\texttt{results/prereg/hpo\_prereg\_v1.1.tar}, holding four files and no others:
the arm configuration \texttt{configs/hpo\_arm.yaml}, the design record
\texttt{results/HPO\_ARM\_SPEC.md}, the search harness
\texttt{scripts/experiments/hpo/asha\_hpo.py} and its QLIKE objective
\texttt{scripts/experiments/hpo/qlike\_loss.py}. % src: results/prereg/hpo_prereg_v1.1.tar (tar -tf lists exactly these four members, verified 2026-08-04)
The archive's SHA-256 digest is published beside it in
\texttt{hpo\_prereg\_v1.1.sha256}, and recomputing the digest over the shipped
archive reproduces the published value
exactly. % src: shasum -a 256 -c results/prereg/hpo_prereg_v1.1.sha256 returns OK (digest 2c9a1c9b0b774ef1f330725a39049a57da14e4e7f9a4423e1335fa5f022f5b59, verified 2026-08-04)
An OpenTimestamps proof file sits alongside the archive, and the directory's own
documentation records that the proof was submitted on 14 July 2026 to three
independent calendar servers and anchored in the Bitcoin blockchain, an
attestation that requires no account and discloses no
identity. % src: results/prereg/README.md (timestamp section: submission date, the three named calendars, Bitcoin anchoring)

That last sentence reports what the documentation records, not a check this
report performed. The OpenTimestamps client is absent from the project's
declared and locked environments, so no attestation was verified while this
appendix was written, and the directory carries a smaller, earlier copy of the
proof beside the current one, which is consistent with the upgrade step the same
documentation describes but is not evidence of what either copy now
contains. % src: results/prereg/ (hpo_prereg_v1.1.tar.ots 4,160 bytes dated 2026-07-23; hpo_prereg_v1.1.tar.ots.bak 1,015 bytes dated 2026-07-14; no opentimestamps entry in pyproject.toml or uv.lock, verified 2026-08-04)
The digest is what this report checked. The anchoring is what a reader can check
independently, by installing the client and running \texttt{ots verify} against
the proof file, using nothing beyond that file and a client.

The instrument's scope should not be inflated either. It covers the tuned arm's
bundle and nothing else. The seven registered analysis records carry
version-control tags but no external stamp, so for those seven the available
claim remains internal consistency rather than temporal priority, exactly as set
out above. What the stamp buys, if its attestation verifies, is priority over
the test evaluation of one arm. It does not buy priority over the author's own
thinking, and it says nothing about the audit protocol as a whole. Marking that
boundary is the point of stating it: a report that repeatedly separates a
checked claim from an asserted one cannot afford to assert this one.

\FloatBarrier
\section{The Benchmark Beside the Corpora It Joins}
\label{app:external:benchmarks}

The primary deliverable of this project is a benchmark, which makes its standing
among the resources it joins a claim to evidence rather than to assume.
Table~\ref{tab:app-benchmarks} sets it beside them property by property, and the
properties are the ones an incremental verdict turns on rather than the ones that
flatter a newcomer.

The comparison set is small because the evidence base is small, and its size was
established rather than guessed. The registered audit of the disclosure-text and
volatility literature recorded a census instead of a sample: three corpora carry
that literature. The models built on the earnings-call side all sit on one of the
three by their own documentation, so counting models rather than corpora would
inflate the denominator without adding evidence. % src: results/tables/kogan_corpus_audit.md (census block: three corpora, MDRM/MAEC/Kogan)
Two of the three were audited end to end by this project and the third could not
be obtained. FETILDA joins the table as the long-document evaluation framework
closest to this study's own construction. FinText is treated below the table
rather than tabulated, for the reason given there.

\begin{table}[p]
  \centering
  \caption[The benchmark beside the corpora it joins]{SP500Vol-Bench beside the
  corpora and evaluation frameworks the disclosure-to-volatility literature was
  built on. Two conventions govern the cells, and the distinction between them is
  the point of the table. \textbf{No} records a property this report
  characterises as absent from the design named, on the reading given in
  Section~\ref{sec:bg-text}.
  \textbf{Unknown} records a property this project did not establish for that
  resource from sources it could check; it is not a claim that the resource lacks
  the property, and it is written rather than inferred. Sizes are what this
  project verified, not what each original reports: the MAEC column gives the
  slice actually audited here, while the other two counted columns give
  whole-corpus totals. The table compares documented construction and evaluation
  properties. It is not a ranking of research quality, and nothing in it is a
  test statistic.}
  \label{tab:app-benchmarks}
  \singlespacing\footnotesize
  \setlength{\tabcolsep}{2pt}
  % Narrow measures are set ragged right so that a 1.85 cm column does not stretch
  % its inter-word spacing; \arraybackslash restores \\ as the row terminator.
  \begin{tabular}{@{}>{\raggedright\arraybackslash}p{2.45cm}>{\raggedright\arraybackslash}p{2.85cm}>{\raggedright\arraybackslash}p{2.5cm}>{\raggedright\arraybackslash}p{1.85cm}>{\raggedright\arraybackslash}p{2.35cm}>{\raggedright\arraybackslash}p{2.15cm}@{}}
    \toprule
    Property & SP500Vol-Bench (10-K/10-Q/8-K, 2010--2025) & Kogan 10-K corpus (1996--2006) & FETILDA (10-K) & MAEC (earnings calls) & MDRM (earnings calls) \\
    \midrule
    Size, as verified here
      & 144{,}129 filings; 427{,}429 modelled rows; 803 filing registrants
      & 30{,}474 filings; 7{,}207 registrants
      & unknown
      & audited slice: 672 test rows; 143 call dates; 461 entities
      & unknown; corpus not obtainable \\ % src: release/accession_index.csv (144,129 filings; 803 filing registrants, TECHNICAL_SUPPLEMENT.md Table S1 block b); frozen 03_dataset.tex Table 1 (427,429 modelled rows); results/tables/kogan_corpus_audit.md (30,474 filings, 7,207 unique CIKs); results/tables/maec_audit.md (672 test rows, 143 call dates, 461 entities per horizon)
    \addlinespace[2pt]
    Label and horizon
      & forward realised volatility from daily closes, modelled in logs; $h = 5/10/20$ trading days
      & log volatility over the 12 months after the filing
      & 10-K forward volatility, the Kogan task
      & post-call realised volatility; $h = 3/7/15/30$ days
      & post-call realised volatility \\ % src: frozen 04_methods.tex (horizons h in {5,10,20}); results/tables/kogan_corpus_audit.md (prereg block: logvol.+12 as the target, logvol.-12 as the control); results/tables/maec_audit.md (HAC lag block: h = 3/7/15/30)
    \addlinespace[2pt]
    Survivorship-free point-in-time universe
      & yes: 914 membership intervals, 411 of them closing inside the sample
      & unknown
      & unknown
      & unknown
      & unknown \\ % src: crsp_restricted/membership_intervals.csv (914 intervals; 411 closing inside the sample)
    \addlinespace[2pt]
    Entity linking audited
      & yes: 30{,}100 point-in-time link windows; 31 successor overrides on 23 firms checked against EDGAR
      & unknown
      & unknown
      & unknown
      & unknown \\ % src: crsp_restricted/cik_links_pit.csv (30,100 link windows); release/DATA_CARD.md lines 45-46 (31 override windows covering 23 firms)
    \addlinespace[2pt]
    Recalibrated price reference
      & yes: validation-fitted and frozen, in log space
      & no
      & no
      & no
      & no \\
    \addlinespace[2pt]
    Firm-identity control
      & yes: identity-augmented reference rung, anonymisation, matched-firm swap
      & no
      & no
      & no
      & no \\
    \addlinespace[2pt]
    Inference clustered on the unit sharing the shock
      & yes: day-clustered DM is the test of record
      & no: filings treated as independent
      & no
      & no
      & no \\
    \addlinespace[2pt]
    Power stated as a minimum detectable effect
      & yes: per-cell MDE beside every null; median 0.82 per cent over 69 cells
      & unknown
      & unknown
      & unknown
      & unknown \\ % src: results/tables/signal_injection_power.md (per-cell MDE; median 0.82 per cent, IQR 0.37--1.27). The four non-project cells are Unknown, not No: this project did not establish what any of the four resources reports about its own power, and the caption reserves No for a property Section 2.1.2 characterises as absent.
    \addlinespace[2pt]
    Availability and licence
      & public core released; CRSP market-data rows and the two CRSP-derived crosswalks withheld; public-price variant covers 72.6/89.3/95.2 per cent of train\slash validation\slash test rows
      & public download from the authors' site; the fetch pipeline ships here, the data do not
      & unknown
      & obtainable and used as published; not redistributed here
      & not obtainable: bundled text and audio, no licence, no longer redistributed \\ % src: results/tables/public_variant_cascade.md (72.6/89.3/95.2 per cent coverage); results/tables/kogan_corpus_audit.md (G-K0 provenance: authors' site, data not redistributed; census block: MDRM status)
    \bottomrule
  \end{tabular}
\end{table}

Read across, the table does not say that SP500Vol-Bench is the larger resource,
and on two dimensions it is the smaller. On filer breadth the Kogan corpus spans
7{,}207 registrants against this benchmark's 803, its universe being the
population of 10-K filers where this one is confined to S\&P~500 membership. % src: results/tables/kogan_corpus_audit.md (7,207 unique CIKs); release/accession_index.csv (803 filing registrants, TECHNICAL_SUPPLEMENT.md Table S1 block b)
That confinement was a decision with a price, and Section~\ref{sec:data-sources}
gives the reason. Realised volatility computed from daily closes is noisiest on
thinly traded stocks, so an index restriction buys a label a day-clustered test
can carry. What it forfeits is whatever the small-firm end of the market would
have shown. On modality the benchmark is poorer again, since both earnings-call
resources align speech to text and a filing has no speech to align.

What the table does show is an asymmetry of documentation on exactly the
properties that decide an incremental verdict. Three of them this report
characterises directly for the incumbent designs, and on its reading no
incumbent imposes any: a recalibrated price reference, a firm-identity control,
and inference clustered on the unit that shares the shock. % src: frozen writing/dissertation/supporting/sections/02_related.tex (none of the closest designs imposes a recalibrated price reference, a firm-identity control, or day-clustered inference); results/tables/maec_audit.md (prereg section 8: call-date-clustered inference never provided in the literature)
On point-in-time construction, on entity linking, and on stated power for three
of the four resources, the entry is \emph{unknown}. It is written rather than
filled in. The pattern of unknowns is itself the finding worth stating: these are
properties a reader of the incumbent resources cannot check, and being checkable
is the condition a benchmark contribution is supposed to supply.

The claim the comparison supports is therefore narrow, and it is the claim the
study makes. To the author's knowledge no prior forecast-level evaluation of this
question jointly imposes a survivorship-free point-in-time panel, a recalibrated
reference interval closed by a maximal price pool, day-clustered inference with
Holm correction and placebo gating, and a model spectrum running from word lists
to a prompted large language model. That is a claim about a protocol and the
fixtures released with it, not about corpus size. What makes it checkable is that
the protocol was carried twice to data assembled by others rather than asserted
at home. On the earnings-call corpus the published-convention gain is reproduced
and then absorbed (Section~\ref{app:res-maec}). On the founding study's own 10-K
corpus the published positive is reproduced at the same sign and order of
magnitude, and the cascade then re-prices it
(Section~\ref{sec:val-kogan}). % src: results/tables/maec_audit.md; results/tables/kogan_corpus_audit.md
A protocol that travels to its predecessors and returns the same verdict is
better positioning evidence than a table, and the table is offered as the map
rather than as the argument.

\textbf{FinText, and why it has no column.} FinText forecasts realised volatility
from word embeddings trained on financial newswire and benchmarks against
HAR-family models, which makes it the one design in the survey that cannot be
charged with a weak price baseline \cite{rahimikia2021fintext}. It has no column
because it is a method and a set of embeddings rather than a released volatility
benchmark. Most of the construction properties compared above would have no entry
for it, so the column would be a run of unknowns. Its lesson is the one the table
is built around: a textual claim is won or lost on the baseline, not on the
corpus.

\FloatBarrier
\section{Neural Price Models, and Why the Reference Here Is Linear}
\label{app:external:neural}

Section~\ref{sec:bg-price} takes the price-only reference as far as the HAR
family and its parametric neighbours and stops. A reader is entitled to ask what
lies further out, because every challenger in this study is neural while the
reference each challenger is measured against is not. This section answers in
the only way the evidence permits: it reports what the literature states, and it
is explicit that this project fitted no neural price model and therefore cannot
adjudicate the question itself.

\textbf{What the literature reports.} Recurrent architectures have been applied
to realised-volatility forecasting directly. Bucci \cite{bucci2020neural}
compares feed-forward and recurrent networks, among them long short-term memory
and nonlinear autoregressive specifications, and reports the recurrent models
outperforming the econometric methods in that comparison, with the variants
built to capture long-range dependence doing best. Kim and Won
\cite{kimwon2018lstmgarch} take the hybrid route, feeding the fitted parameters
of several GARCH-type models into a long short-term memory network and reporting
the hybrid ahead of any single GARCH specification on a stock index. Closest to
this project's own question is the comparison of Christensen, Siggaard and
Veliyev \cite{christensen2023mlvol}. They set regularised regressions, regression
trees and neural networks against several HAR specifications, on realised
variance for the constituents of the Dow Jones Industrial Average, and report
machine learning beating the HAR lineage even when the only predictors are the
daily, weekly and monthly realised-variance lags HAR itself uses, with the margin
widening at longer horizons. Their attribution matters as much as the verdict:
the gain comes from greater fitted persistence, which approximates the long
memory of realised variance more closely than a three-lag linear form does.

\textbf{What this study can and cannot say.} It cannot say that HAR beats a
neural price model, because it never fitted one, and no statement anywhere in
this report should be read as a comparison against a neural price baseline that
was tested. Nor can it claim that its reference is the best price forecast
obtainable on this panel. Section~\ref{sec:bg-price} sets the standard that a
credible audit allows the price side its best available representative. This is
the one direction in which the study meets that standard only in part. The rungs
of the reference ladder are recalibration, a firm-identity term and a pool of
linear and parametric specifications, and not one rung is a learned nonlinear
price model.

\textbf{Why a linear reference was nevertheless the right one here.} Three
reasons, in decreasing order of force. The first is protocol parity. HAR-RV is
estimated by ordinary least squares, as is every combination the incremental
instrument fits, so a comparison against it isolates the information source. A
nonlinear learner on the reference side facing a nonlinear learner on the
challenger side would instead confound the information source with estimator
capacity, which is the confusion this design exists to prevent. The
second is that the question is incremental rather than absolute. What the
reference must be is well calibrated and identity-aware, because an uncorrected
reference donates its own miscalibration to the challenger as credit, and the
ladder pays for those two properties explicitly rather than for raw accuracy.
The third is direction, and it is measured rather than argued. Two executed
tests say what a stronger price reference does to the text verdict. The first is
the VIX-augmented re-test of Table~\ref{tab:app-ladder}, which leaves 19 of the 38
re-tested cells standing; the second is the frontier completion of
Section~\ref{app:res-frontier}.
% src: results/tables/vix_control.md (19 of the 38, seed 2026, Holm in 40)
Adding a VIX-augmented HAR and a
measurement-error-corrected HAR to the pool cuts the Holm survivor count from 8
of 69 cells to 3 of 69, on the archived single-seed basis both readings share. % src: results/tables/pool_frontier_audit.md (Holm survivors 8 of 69 under pool5, 3 of 69 under pool7); results/tables/pool_frontier_cascade.csv (genuine_holm, per-cell)
Strengthening the price side has therefore been shown, on this panel, to deepen
the null rather than to relieve it.

\textbf{What the gap costs, and why it cannot now be closed.} An omitted stronger
reference threatens a positive finding, and the finding on text here is a bounded
near-null, so the direction of the threat runs against the residual this study
does report rather than in its favour. Nothing on this evidence suggests a neural
rung would run the other way, though that expectation is an inference from the
frontier test and not a result. Closing the gap would require a new fit, and the
experimental record is frozen: the training fleet is decommissioned, so no rung
can now be added and re-scored under the protocol the rest of the ladder ran.
The gap is recorded as a limitation of scope, of the same class as the arms
Section~\ref{sec:conc-limitations} already lists rather than a new one.
